\documentclass[12pt]{article}
\usepackage{amsmath,amssymb,amsthm}
\usepackage{fullpage}
\usepackage{algorithm}
\usepackage{algorithmic}

\newtheorem{theorem}{Theorem}
\newtheorem{lemma}{Lemma}
\newtheorem{definition}{Definition}

\begin{document}

\section*{Algorithm Name}
Online Gradient Descent for Strongly Convex Online Functions (OGD-SC)

\section*{Mathematical Formulation}
We consider an online sequence of loss functions $\{f_t:\mathcal{W}\to\mathbb{R}\}_{t=1}^T$ revealed in rounds $t=1,2,\dots,T$.
- Decision set: $\mathcal{W}\subseteq \mathbb{R}^d$ is a nonempty closed convex set.
- At round $t$, the learner picks $w_t\in \mathcal{W}$, then observes $f_t$ and incurs loss $f_t(w_t)$. A subgradient (or gradient) $g_t\in \partial f_t(w_t)$ is revealed and used to update.
- Update rule (projected OGD):
\[
w_{t+1} \;=\; \Pi_{\mathcal{W}}\!\left(w_t - \eta_t \, g_t\right),
\]
where $\eta_t>0$ is a step size and $\Pi_{\mathcal{W}}$ denotes Euclidean projection onto $\mathcal{W}$:
\[
\Pi_{\mathcal{W}}(x) \;\triangleq\; \arg\min_{w\in \mathcal{W}} \|w-x\|_2.
\]
- Step-size schedule for strongly convex losses:
\[
\eta_t \;=\; \frac{1}{\lambda t}, \quad t\ge 1,
\]
where $\lambda>0$ is the strong convexity parameter.

Performance is measured by the regret after $T$ rounds:
\[
R_T \;\triangleq\; \sum_{t=1}^T f_t(w_t) \;-\; \min_{w\in \mathcal{W}} \sum_{t=1}^T f_t(w).
\]

\section*{Pseudocode}
\begin{algorithm}[H]
\caption{OGD-SC: Online Gradient Descent for Strongly Convex Losses}
\begin{algorithmic}[1]
\STATE \textbf{Input:} Strong convexity parameter $\lambda>0$, horizon $T$, decision set $\mathcal{W}\subseteq\mathbb{R}^d$
\STATE \textbf{Initialize:} Pick any $w_1\in \mathcal{W}$
\FOR{$t=1$ to $T$}
    \STATE Receive (or evaluate) loss $f_t:\mathcal{W}\to\mathbb{R}$
    \STATE Compute a subgradient $g_t\in \partial f_t(w_t)$
    \STATE Set step size $\eta_t = 1/(\lambda t)$
    \STATE Update $w_{t+1} = \Pi_{\mathcal{W}}\!\left(w_t - \eta_t g_t\right)$
\ENDFOR
\STATE \textbf{Output:} Sequence $\{w_t\}_{t=1}^T$ (or an average if desired)
\end{algorithmic}
\end{algorithm}

\section*{Dry Run Example}
We apply the basic (unconstrained) gradient update to the single function $f(w) = (w-3)^2$ in one dimension with $w_1=0$, constant step size $\eta=0.1$, for $3$ iterations. The gradient is $g(w)=\nabla f(w)=2(w-3)$.
- Iteration $t=1$:
  - Current $w_1=0$
  - Gradient: $g_1 = 2(0-3) = -6$
  - Update: $w_2 = w_1 - \eta g_1 = 0 - 0.1(-6) = 0.6$
  - Objective: $f(w_1) = (0-3)^2 = 9$
- Iteration $t=2$:
  - Current $w_2=0.6$
  - Gradient: $g_2 = 2(0.6-3) = -4.8$
  - Update: $w_3 = 0.6 - 0.1(-4.8) = 1.08$
  - Objective: $f(w_2) = (0.6-3)^2 = 5.76$
- Iteration $t=3$:
  - Current $w_3=1.08$
  - Gradient: $g_3 = 2(1.08-3) = -3.84$
  - Update: $w_4 = 1.08 - 0.1(-3.84) = 1.464$
  - Objective: $f(w_3) = (1.08-3)^2 = 3.6864$

\section*{Assumptions}
We assume the following throughout the analysis:
- A1 (Convex domain): $\mathcal{W}\subseteq \mathbb{R}^d$ is a nonempty closed convex set.
- A2 (Strong convexity): Each $f_t:\mathcal{W}\to\mathbb{R}$ is $\lambda$-strongly convex with respect to the Euclidean norm, i.e., for all $u,v\in \mathcal{W}$ and any $g_t(v)\in \partial f_t(v)$,
\[
f_t(u) \;\ge\; f_t(v) + \langle g_t(v), u-v\rangle + \frac{\lambda}{2}\|u-v\|^2.
\]
- A3 (Bounded subgradients): There exists $G>0$ such that for all $t$ and all $w\in \mathcal{W}$, any subgradient $g\in \partial f_t(w)$ satisfies $\|g\|\le G$.
- A4 (Step size): $\eta_t = \frac{1}{\lambda t}$.

\section*{Main Convergence Theorem}
\begin{theorem}[Logarithmic Regret for Strongly Convex Online Functions]
Under Assumptions A1--A4, the regret of OGD-SC satisfies
\[
R_T \;\le\; \frac{G^2}{2\lambda}\left(1 + \ln T\right), \quad \text{for all } T\ge 1.
\]
\end{theorem}

\section*{Theoretical Properties}
- Stability: The update satisfies $\|w_{t+1}-w_t\|\le \eta_t \|g_t\|\le \eta_t G=\frac{G}{\lambda t}$, which vanishes as $t$ grows, implying stability of the iterates.
- Hyperparameter sensitivity: The schedule $\eta_t=\frac{1}{\alpha t}$ yields $R_T\le \frac{G^2}{2\alpha}(1+\ln T)$ provided $\alpha\le \lambda$ (see Special Cases), with the best constant at $\alpha=\lambda$.
- Comparison to convex (not strongly convex) case: For merely convex Lipschitz losses, OGD achieves $R_T=O(\sqrt{T})$ with an appropriate constant step size or $\eta_t\propto 1/\sqrt{t}$; strong convexity improves this to $O(\log T)$.

\section*{Discussion}
The algorithm does not target convergence to a fixed point; instead, it minimizes regret against the best fixed comparator in hindsight. Strong convexity sharpens the per-round stability and enables a telescoping argument that removes dependence on the domain diameter, yielding a logarithmic regret bound. The bound holds against an arbitrary (possibly adversarial) sequence of strongly convex losses satisfying the gradient bound.

\section*{Preliminaries (Definitions and Lemmas)}
\begin{definition}[Regret]
For a sequence $\{f_t\}_{t=1}^T$ and actions $\{w_t\}_{t=1}^T$, the regret is
\[
R_T \;=\; \sum_{t=1}^T f_t(w_t) - \min_{w\in \mathcal{W}} \sum_{t=1}^T f_t(w).
\]
Let $w^\star\in\arg\min_{w\in \mathcal{W}} \sum_{t=1}^T f_t(w)$ be any minimizer.
\end{definition}

\begin{lemma}[Euclidean Projection Inequality]
Let $\Pi_{\mathcal{W}}$ be the Euclidean projection onto a nonempty closed convex set $\mathcal{W}$. For any $x\in\mathbb{R}^d$ and any $y\in \mathcal{W}$,
\[
\|\Pi_{\mathcal{W}}(x)-y\|^2 \;\le\; \|x-y\|^2 - \|x-\Pi_{\mathcal{W}}(x)\|^2 \;\le\; \|x-y\|^2.
\]
\end{lemma}
\begin{proof}
Let $p=\Pi_{\mathcal{W}}(x)$. By the first-order optimality condition of the projection onto a convex set, we have $\langle x-p, y-p\rangle \le 0$ for all $y\in \mathcal{W}$.
\[
\|x-y\|^2 = \|x-p+p-y\|^2 = \|x-p\|^2 + \|p-y\|^2 + 2\langle x-p, p-y\rangle. \quad \text{[Step 1: expand square]}
\]
Using $\langle x-p, y-p\rangle \le 0$ gives $\langle x-p, p-y\rangle \ge 0$. Hence
\[
\|x-y\|^2 \;\ge\; \|x-p\|^2 + \|p-y\|^2. \quad \text{[Step 2: drop nonnegative cross term]}
\]
Rearranging yields $\|p-y\|^2 \le \|x-y\|^2 - \|x-p\|^2 \le \|x-y\|^2$. \quad \text{[Step 3: algebraic rearrangement]}
\end{proof}

\begin{lemma}[One-Step Progress Inequality for OGD]\label{lem:one-step}
Let $w_{t+1} = \Pi_{\mathcal{W}}(w_t - \eta_t g_t)$, with any $g_t\in\mathbb{R}^d$ and $\eta_t>0$. Then for any $w^\star\in \mathcal{W}$,
\[
\langle g_t, w_t - w^\star\rangle \;\le\; \frac{\|w_t - w^\star\|^2 - \|w_{t+1}-w^\star\|^2}{2\eta_t} \;+\; \frac{\eta_t}{2}\,\|g_t\|^2.
\]
\end{lemma}
\begin{proof}
Let $x_t = w_t - \eta_t g_t$. Then by Lemma 1 with $x=x_t$, $y=w^\star$, and $p=w_{t+1}=\Pi_{\mathcal{W}}(x_t)$,
\[
\|w_{t+1}-w^\star\|^2 \;\le\; \|x_t - w^\star\|^2. \quad \text{[Step 4: apply projection inequality]}
\]
Expand $\|x_t - w^\star\|^2$:
\[
\|x_t - w^\star\|^2 = \|w_t - \eta_t g_t - w^\star\|^2 = \|w_t - w^\star\|^2 - 2\eta_t \langle g_t, w_t - w^\star\rangle + \eta_t^2 \|g_t\|^2. \quad \text{[Step 5: expand square]}
\]
Combine the two displays:
\[
\|w_{t+1}-w^\star\|^2 \;\le\; \|w_t - w^\star\|^2 - 2\eta_t \langle g_t, w_t - w^\star\rangle + \eta_t^2 \|g_t\|^2. \quad \text{[Step 6: transitivity]}
\]
Rearrange to isolate $\langle g_t, w_t - w^\star\rangle$:
\[
2\eta_t \langle g_t, w_t - w^\star\rangle \;\le\; \|w_t - w^\star\|^2 - \|w_{t+1}-w^\star\|^2 + \eta_t^2 \|g_t\|^2. \quad \text{[Step 7: algebra]}
\]
Divide by $2\eta_t>0$ to conclude:
\[
\langle g_t, w_t - w^\star\rangle \;\le\; \frac{\|w_t - w^\star\|^2 - \|w_{t+1}-w^\star\|^2}{2\eta_t} + \frac{\eta_t}{2}\|g_t\|^2. \quad \text{[Step 8: divide by }2\eta_t]
\]
\end{proof}

\begin{lemma}[Strong Convexity Upper Bound]\label{lem:sc}
If $f_t$ is $\lambda$-strongly convex and $g_t\in \partial f_t(w_t)$, then for any $w^\star\in \mathcal{W}$,
\[
f_t(w_t) - f_t(w^\star) \;\le\; \langle g_t, w_t - w^\star\rangle \;-\; \frac{\lambda}{2}\,\|w_t - w^\star\|^2.
\]
\end{lemma}
\begin{proof}
By strong convexity with $u=w^\star$ and $v=w_t$,
\[
f_t(w^\star) \;\ge\; f_t(w_t) + \langle g_t, w^\star - w_t\rangle + \frac{\lambda}{2}\|w^\star - w_t\|^2. \quad \text{[Step 9: apply definition]}
\]
Rearrange:
\[
f_t(w_t) - f_t(w^\star) \;\le\; \langle g_t, w_t - w^\star\rangle - \frac{\lambda}{2}\|w_t - w^\star\|^2. \quad \text{[Step 10: move terms, use symmetry of norm]}
\]
\end{proof}

\begin{lemma}[Harmonic Series Bound]\label{lem:harm}
For all integers $T\ge 1$, $\sum_{t=1}^T \frac{1}{t} \le 1 + \ln T$.
\end{lemma}
\begin{proof}
For $t\ge 2$, $\int_{t-1}^{t}\frac{1}{x}\,dx = \ln t - \ln(t-1) \le \frac{1}{t-1}$. Hence
\[
\sum_{t=2}^T \frac{1}{t} \;\le\; \int_{1}^{T} \frac{1}{x}\,dx \;=\; \ln T. \quad \text{[Step 11: integral comparison]}
\]
Add $\frac{1}{1}=1$ to both sides to conclude the claim. \quad \text{[Step 12: add $t=1$ term]}
\end{proof}

\section*{Main Proof (Step-by-Step Derivation)}
\begin{proof}[Proof of Theorem 1]
Fix $w^\star\in \arg\min_{w\in \mathcal{W}} \sum_{t=1}^T f_t(w)$. For each $t$, apply Lemma \ref{lem:sc} and then Lemma \ref{lem:one-step}:
\begin{align*}
f_t(w_t) - f_t(w^\star) 
&\le \langle g_t, w_t - w^\star\rangle - \frac{\lambda}{2}\|w_t - w^\star\|^2 \quad \text{[Step 13: Lemma \ref{lem:sc}]}\\
&\le \frac{\|w_t - w^\star\|^2 - \|w_{t+1}-w^\star\|^2}{2\eta_t} + \frac{\eta_t}{2}\|g_t\|^2 - \frac{\lambda}{2}\|w_t - w^\star\|^2. \quad \text{[Step 14: Lemma \ref{lem:one-step}]}
\end{align*}
Sum the inequality over $t=1$ to $T$:
\begin{align*}
\sum_{t=1}^T \bigl(f_t(w_t) - f_t(w^\star)\bigr)
&\le \sum_{t=1}^T \left[\frac{\|w_t - w^\star\|^2 - \|w_{t+1}-w^\star\|^2}{2\eta_t} - \frac{\lambda}{2}\|w_t - w^\star\|^2 \right] + \frac{1}{2}\sum_{t=1}^T \eta_t \|g_t\|^2. \quad \text{[Step 15: sum both sides]}
\end{align*}
Introduce $u_t \triangleq \|w_t - w^\star\|^2$ and $c_t \triangleq \frac{1}{\eta_t}$. Then the bracketed sum equals
\begin{align*}
\frac{1}{2}\sum_{t=1}^T \bigl(c_t(u_t - u_{t+1}) - \lambda u_t\bigr)
&= \frac{1}{2}\left[ c_1 u_1 - c_T u_{T+1} + \sum_{t=2}^T (c_t - c_{t-1})u_t - \lambda \sum_{t=1}^T u_t \right] \quad \text{[Step 16: telescoping identity]}\\
&= \frac{1}{2}\left[ (c_1-\lambda)u_1 - c_T u_{T+1} + \sum_{t=2}^T \bigl((c_t - c_{t-1}) - \lambda\bigr)u_t \right]. \quad \text{[Step 17: regroup terms]}
\end{align*}
Now specialize to $\eta_t = \frac{1}{\lambda t}$ so that $c_t=\lambda t$ is affine in $t$:
\[
c_1-\lambda = \lambda - \lambda = 0, \quad c_t - c_{t-1} - \lambda = \lambda - \lambda = 0 \text{ for all } t\ge 2. \quad \text{[Step 18: compute differences]}
\]
Therefore the entire bracketed sum reduces to
\[
\frac{1}{2}\left[ -c_T u_{T+1}\right] \;\le\; 0, \quad \text{since } c_T u_{T+1}\ge 0. \quad \text{[Step 19: nonnegativity implies upper bound by 0]}
\]
Discarding this nonpositive term yields
\[
\sum_{t=1}^T \bigl(f_t(w_t) - f_t(w^\star)\bigr) \;\le\; \frac{1}{2}\sum_{t=1}^T \eta_t \|g_t\|^2. \quad \text{[Step 20: drop nonpositive contribution]}
\]
Using Assumption A3, $\|g_t\|\le G$ for all $t$, and $\eta_t = \frac{1}{\lambda t}$,
\[
\sum_{t=1}^T \bigl(f_t(w_t) - f_t(w^\star)\bigr) \;\le\; \frac{1}{2}\sum_{t=1}^T \frac{G^2}{\lambda t} \;=\; \frac{G^2}{2\lambda}\sum_{t=1}^T \frac{1}{t}. \quad \text{[Step 21: bound gradients and plug step sizes]}
\]
Apply Lemma \ref{lem:harm} to bound the harmonic series:
\[
\sum_{t=1}^T \bigl(f_t(w_t) - f_t(w^\star)\bigr) \;\le\; \frac{G^2}{2\lambda}\left(1 + \ln T\right). \quad \text{[Step 22: Lemma \ref{lem:harm}]}
\]
Since the left-hand side equals the regret $R_T$ by definition, the theorem follows. \quad \text{[Step 23: identify regret]}
\end{proof}

\section*{Rate Derivation (With Constants)}
With $\eta_t = \frac{1}{\lambda t}$ and $\|g_t\|\le G$,
\[
R_T \;\le\; \frac{1}{2}\sum_{t=1}^T \frac{G^2}{\lambda t} 
\;\le\; \frac{G^2}{2\lambda}\left(1+\ln T\right),
\]
where:
- $G$ is the uniform bound on subgradient norms (A3),
- $\lambda$ is the strong convexity parameter (A2),
- The constant $1$ in $1+\ln T$ comes from the $t=1$ harmonic term.

\section*{Special Cases}
- Unconstrained case $(\mathcal{W}=\mathbb{R}^d)$: The projection becomes identity, i.e., $w_{t+1}=w_t - \eta_t g_t$. Lemma \ref{lem:one-step} holds with equality at Step 4, and the same regret bound follows.
- Misspecified step size $\eta_t = \frac{1}{\alpha t}$ with $\alpha \le \lambda$: Then $c_t=\alpha t$ and $c_t-c_{t-1}-\lambda = \alpha - \lambda \le 0$. In Steps 16--20, the bracketed sum is $\le 0$ and can be dropped, yielding
\[
R_T \;\le\; \frac{G^2}{2\alpha}\sum_{t=1}^T \frac{1}{t} \;\le\; \frac{G^2}{2\alpha}\left(1+\ln T\right),
\]
which is largest when $\alpha$ is much smaller than $\lambda$; the best constant occurs at $\alpha=\lambda$.
- Overaggressive schedule $\eta_t = \frac{1}{\alpha t}$ with $\alpha>\lambda$: Then $c_t-c_{t-1}-\lambda=\alpha-\lambda>0$, and the bracketed sum in Step 17 contributes a positive term that requires an a priori bound on $\|w_t-w^\star\|^2$ (e.g., via bounded $\mathcal{W}$). Without such a bound, the logarithmic guarantee may fail; thus choosing $\alpha\le \lambda$ is recommended for diameter-free regret.

\end{document}